\section{HTML}\label{sec:HTML}

\begin{frame}
    \frametitle{Hypertext Markup Language (HTML)}
    
    HTML es un lenguaje para armar documentos. ``Hypertext'' se 
    refiere a los links que conectan los distintos sitios y documentos. 
    ``Markup'' se refiere a que se usa una sintaxis particular para 
    marcar cómo se debería estructurar el contenido.

    La sintaxis de HTML consiste en el uso de \emph{tags}. Hay alrededor de 
    140 de estas tags, pero lo normal es usar un subconjunto bastante 
    más chico de ellas. 
    
    Las tags tienen el formato \emph{$<$tag propiedades$>$Contenido$</$tag$>$} o \\
    \emph{$<$tag propiedades$/>$} para elementos sin contenido. Dicho contenido 
    podría ser texto o incluso otras tags. Veamos cuáles son las tags más importantes.

\end{frame}

\begin{frame}[fragile]
    \frametitle{Estructura principal}
    
    Las tags más importante que cualquier sitio debe tener son \textbf{html}, 
    \textbf{head} y \textbf{body}. La primera es la que define los límites del documento, 
    sin esa tag no existe nada. La segunda es para definir información relevante 
    al sitio pero que no va a ser visible para el usuario, eso incluye metadatos 
    y un poco de información extra para que el navegador sepa como interpretar el 
    sitio. Por último, \textbf{body} es donde va el contenido del sitio, básicamente todo.

    Cualquier sitio normal sigue el siguiente esquema:

    \begin{lstlisting}
<html>
    <head>
        Metadatos, estilo, etc...
    </head>
    <body>
        Contenido
    </body>
</html>
    \end{lstlisting}

\end{frame}

\begin{frame}
    \frametitle{Contenidos de \emph{head}, metadatos}
    
    \begin{itemize}
        \item \textbf{title} se usa para elegir el título de la página. Dicho 
              título es el que se ve en la pestaña correspondiente del navegador.
        \item \textbf{link} define relaciones con recursos externos al documento. 
              Se usa principalmente para elegir la hoja de estilos (CSS) pero 
              también otras cosas como el ícono del sitio (suele aparecer al 
              lado del título). En el caso más general usa la propiedad 
              \emph{href} para definir la dirección del recurso y \emph{rel} para 
              especificar cual es dicha relación.
        \item \textbf{style} permite agregar estilo directamente ahí pero no es muy 
              común hacer eso.
        \item \textbf{meta} sirve para agregar cualquier otro tipo de información 
              que no se pueda representar con las tags anteriores.
    \end{itemize}

\end{frame}

\begin{frame}
    \frametitle{Estructura principal pero menos principal}

    HTML viene con los elementos \textbf{header}, \textbf{main} y \textbf{footer} para 
    organizar el sitio en esas tres partes. También incluye cosas como \textbf{aside} 
    para cosas que justamente deberían ir a un costado, \textbf{article} para artículos, 
    etc. (todas bastante auto explicativas).
    
    Sin embargo, suele pasar (y es discutible si es algo bueno o malo) que el único 
    elemento que se termina usando para organizar el sitio es \textbf{div}. El elemento 
    justamente es un contenedor genérico que no representa nada por si mismo.

    Todos estos elementos tienen formato ``block'', que, ahora mismo y sin entrar 
    en detalles, significa que ocupan el ancho entero del sitio. También implica otras 
    cosas pero todo eso va a importar más cuando nos pongamos con los estilos.

\end{frame}

\begin{frame}
    \frametitle{Texto}

    Los elementos \textbf{h$N$}, donde la $N$ va de 1 a 6, son encabezados. La importancia 
    va de menor a mayor, donde \textbf{h1} es el título principal del documento (y debería 
    haber solo uno) y \textbf{h6} es un título bastante menor.

    El elemento \textbf{p} representa un párrafo, aunque se puede meter todo el texto que 
    quieran ahí adentro.

    Esos son los principales elementos que contienen texto y ambos tienen formato block.
    Pero también hay elementos que modifican el texto que está adentro de los \textbf{p} 
    y \textbf{h}.

    \textbf{a} es para agregar links, \textbf{strong} es para marcar importancia, 
    \textbf{em} es para énfasis, etc. Al igual que \textbf{div}, la versión genérica 
    de los modificadores es \textbf{span}. Todos estos elementos tienen formato 
    ``inline'', lo cual por ahora solo implica que van en el medio del contenido.

\end{frame}

\begin{frame}
    \frametitle{Listas y tablas}

    También tenemos listas y tablas.

    \begin{itemize}
        \item \textbf{ul} es para listas no ordenadas (el punteado no tiene números) y 
              \textbf{ol} es para listas ordenadas (el punteado si tiene números (o letras
              (o lo que quieran (se puede cambiar todo)))). Los elementos que van adentro 
              de estas listas deberían ser \textbf{li}.
        \item \textbf{table} es para tablas. Hay un montón de elementos relacionados a las 
              tablas para organizarlas como quieran, pero los más importantes son 
              \textbf{tr} para filas, \textbf{th} para cabezales de columnas y \textbf{td}
              para celdas..
    \end{itemize}

\end{frame}

\begin{frame}
    \frametitle{Imágenes y formularios}

    Para terminar con los elementos más comunes tenemos imágenes y formularios.

    \begin{itemize}
        \item Para usar imágenes está \textbf{img}. La propiedad \emph{src} define la dirección
              de la imagen a mostrar. Tiene unas cuantas propiedades más para cambiar como 
              aparece y tocar temas de accesibilidad.
        \item Los formularios son con \textbf{form} y los campos de dichos formularios 
              son con \textbf{input}. La propiedad \emph{type} define que tipo de 
              campo es.
    \end{itemize}

\end{frame}

\begin{frame}
    \frametitle{Ejercicio - Perfil}

    En este taller cada los ejercicios están divididos en distintos archivos .html 
    y .css. En la carpeta ``ejercicios'' están los archivos a completar y en 
    ``resueltos'' los archivos completos. La idea es que solo miren los resueltos 
    al final para comparar con lo que hicieron o que a lo sumo los abran (en el 
    navegador, no un editor de texto) para ver cómo podrían hacer que se vea el 
    resultado. En los archivos a completar se pueden ver las consignas comentadas.

    Este primer ejercicio es solo de html y se puede resolver usando las tags que 
    vimos hasta ahora. La idea es que nos cuenten un poco sobre ustedes y 
    sobre la relacion que tienen con la carrera. 

\end{frame}